(1) 弦的非中点位置满足 $\mu y_{tt}=Ty_{xx}$，端点 $y(0,t)=y(L,t)=0$，中点 $y$ 连续。对中点点质量应用牛顿定律得 $m y_{tt}(L/2,t)=T[y_x(L/2^+,t)-y_x(L/2^-,t)]$；这里导数跳变来自作用在点质量两侧的横向张力差。
(2) 反对称模在中点位移为零且左右斜率连续，故 $kL/2=n\pi$，$n=1,2,\ldots$。对称模在左半弦取 $A\sin(kx)$，右半弦取 $A\sin[k(L-x)]$；中点跳变方程给 $-m\omega^2A\sin(kL/2)=-2TkA\cos(kL/2)$，即 $\cot(kL/2)=mk/(2\mu)$，并有 $\omega=ck$、$c=\sqrt{T/\mu}$。
(3) 当 $m\to0$，对称条件退化为 $\cos(kL/2)=0$，与反对称条件合并后得普通弦全部 $k_j=j\pi/L$。当 $m\gg\mu L$，最低对称根很小，$\cot(kL/2)\sim2/(kL)$，所以 $k^2\sim4\mu/(mL)$、$\omega^2\sim4T/(mL)$，与两半弦等效总回复刚度 $4T/L$ 一致。
(4) 弦与点质量的总能量为 $E=\frac12\int_0^L[\mu y_t^2+Ty_x^2]dx+\frac12m y_t(L/2,t)^2$。分段对时间求导和分部积分后，中点边界功率项恰被点质量的动能导数抵消，故 $dE/dt=0$；额外质量使低频模变慢，但不是阻尼器。
